\documentclass[12pt]{article}

\usepackage{bbm}
\usepackage{fancyhdr}
\usepackage{titling}
\usepackage{enumitem}
\usepackage{amsmath}
\usepackage{dsfont}
\usepackage{amsfonts}
\usepackage{amssymb}
\usepackage{amsthm}
\usepackage{geometry}
\usepackage{tikz}
\usepackage{algorithm}
\usepackage{algpseudocode}
\geometry{verbose,tmargin=1in,bmargin=1in,lmargin=1in,rmargin=1in}
\renewcommand{\vec}[1]{\mathbf{#1}}
\newcommand{\dec}{\text{dec}}
\newcommand{\divi}{\mathbin{\textnormal{div}}}
\title{Homework 1}
\author{Satvik Prasad}
\date{}
\newtheorem{lemma}{Lemma}[section]
\fancypagestyle{plain}{%
  \fancyhf{}
  \fancyhead[L]{\theauthor}
  \fancyhead[R]{\thedate}
}
\pagestyle{plain}
\makeatletter
\def\@maketitle{%
  \newpage
  \null
  \vskip 1em%
  \begin{center}%
    \let \footnote \thanks
    {\LARGE \@title \par}%
    \vskip 1em%
  \end{center}%
  \par
\vskip 1em}
\makeatother
\usepackage{cmbright}

\newcommand{\mat}[1]{\mathbf{#1}}

% ---------- Probability theory ----------
% Core operators (optional argument-free; use as \E{X}, \Var{X}, \Cov{X}{Y})
\DeclareMathOperator{\EOp}{\mathbb{E}}
\DeclareMathOperator{\VarOp}{Var}
\DeclareMathOperator{\CovOp}{Cov}
\DeclareMathOperator{\CorrOp}{Corr}

\renewcommand{\P}[1]{\mathbb{P}\left(#1\right)}        % \P{A}
\newcommand{\E}[1]{\EOp\left[#1\right]}                % \E{X}
\newcommand{\Var}[1]{\VarOp\left(#1\right)}            % \Var{X}
\newcommand{\Cov}[2]{\CovOp\left(#1,\,#2\right)}       % \Cov{X}{Y}
\newcommand{\Corr}[2]{\CorrOp\left(#1,\,#2\right)}     % \Corr{X}{Y}
\newcommand{\SD}[1]{\operatorname{SD}\left(#1\right)}  % \SD{X}

% Subscripted versions, e.g. \Esub{\theta}{X}, \Psub{\theta}{A}
\newcommand{\Esub}[2]{\EOp_{#1}\left[#2\right]}
\newcommand{\Psub}[2]{\mathbb{P}_{#1}\left(#2\right)}

% Conditioning, indicators, independence
\newcommand{\given}{\,\middle|\,}                      % \E{X \given Y}
\newcommand{\ind}[1]{\mathbbm{1}\left\{#1\right\}}     % \ind{X > 0}
\newcommand{\indep}{\perp\!\!\!\perp}                  % X \indep Y
\newcommand{\iid}{\overset{\text{iid}}{\sim}}          % X_i \iid \Normal(0,1)

% Convergence
\newcommand{\convp}{\xrightarrow{\;p\;}}               % in probability
\newcommand{\convd}{\xrightarrow{\;d\;}}               % in distribution
\newcommand{\convas}{\xrightarrow{\text{a.s.}}}        % almost surely
\newcommand{\convLp}[1]{\xrightarrow{L^{#1}}}          % in L^p

% Probability spaces
\newcommand{\F}{\mathcal{F}}                           % sigma-algebra
\newcommand{\B}{\mathcal{B}}                           % Borel sigma-algebra
\newcommand{\pspace}{(\Omega, \F, \mathbb{P})}

% Common distributions
\newcommand{\Normal}{\mathcal{N}}
\newcommand{\Bern}{\operatorname{Bernoulli}}
\newcommand{\Bin}{\operatorname{Binomial}}
\newcommand{\Pois}{\operatorname{Poisson}}
\newcommand{\Unif}{\operatorname{Uniform}}
\newcommand{\Exp}{\operatorname{Exp}}
\newcommand{\Geom}{\operatorname{Geometric}}
\newcommand{\Gam}{\operatorname{Gamma}}
\newcommand{\Beta}{\operatorname{Beta}}

% Number sets
\newcommand{\R}{\mathbb{R}}
\newcommand{\N}{\mathbb{N}}
\newcommand{\Z}{\mathbb{Z}}


\begin{document}
\subsection{Writing Proofs}
\begin{enumerate}
  \item Since the inner product is linear in the first argument and
    conjugate linear in the second, we have
    \begin{align*}
      S &:= \sum_{i = 1} ^{n} \sum_{j = 1}^{n} \langle v_i, v_j
      \rangle = \sum_{j = 1}^{n} \sum_{i = 1}^{n} \langle v_i,
      v_j\rangle = \sum_{j = 1}^{n} \left\langle \sum_{i = 1}v_i, v_j
      \right\rangle \\
      &= \left\langle \sum_{i = 1}^{n} v_i, \sum_{j = 1}^{n}
      v_j\right\rangle = \left|\left| \sum_{i =
      1}^{n}v_i\right|\right|^2 \geq 0,
    \end{align*}
    and hence $S \geq 0$ as required.

  \item Starting from the RHS, we have
    \begin{align*}
      \text{RHS} =  \dfrac{||u + v||^2 - ||u - v||^2}{4} =
      \dfrac{\langle u + v, u + v\rangle - \langle u - v, u - v\rangle}{4}.
    \end{align*}
    By linearity in the first argument and conjugate linearity in the
    second, we have
    \begin{align*}
      \text{RHS}  &= \dfrac{\langle u, u \rangle + \langle v, u
        \rangle + \langle u, v \rangle + \langle v, v \rangle -
        (\langle u, u \rangle + \langle v, v \rangle - \langle v, u
      \rangle - \langle u, v \rangle )}{4} \\
      &= \dfrac{2 \langle v, u \rangle + 2 \langle u, v\rangle}{4}.
    \end{align*}
    Since the field is over $\mathbb R, $ $\bar \lambda = \lambda,$
    meaning the inner product is symmetric, yielding $$
    \text{RHS} = \dfrac{4 \langle u, v \rangle}4{} = \langle u, v
    \rangle = \text{LHS}
    $$
    as required.

  \item
    \begin{enumerate}
      \item Starting from the LHS, we have
        \begin{align*}
          \dfrac{d}{dt} \langle f(t), g(t) \rangle = \dfrac{d}{dt}
          \sum_{i = 1}^{n} f_i(t) g_i(t).
        \end{align*}
        Using the product rule, we get
        $$
        \dfrac{d}{dt} \langle f(t), g(t) \rangle = \sum_{i = 1}^{n}
        f_i(t) g_i'(t) + f_i'(t) g_i(t) = \sum_{i = 1}^{n} f_i(t)
        g_i'(t) + \sum_{i = 1}^{n} f_i'(t) g(t).
        $$
        Therefore, we conclude $$
        \dfrac{d}{dt} \langle f(t), g(t)\rangle = \langle f'(t),
        g(t)\rangle + \langle f(t),g'(t)\rangle
        $$
        as required.

      \item Suppose we have a $c > 0$ for which $||f(t)|| = c$  for
        all $t \in \mathbb R.$ Then, $\langle f(t), f(t) \rangle =
        c^2,$ and differentiating yields
        $$
        \langle f'(t), f(t)\rangle + \langle f(t), f'(t)\rangle = 0.
        $$
        Since this is over the field $\mathbb R$, the inner product
        is symmetric, and we can therefore conclude
        $$
        \langle f'(t), f(t)\rangle = 0
        $$
        for all $t \in \mathbb R.$

      \item Tangent vector to curve is orthogonal to displacement.
    \end{enumerate}

  \item Let us prove the forward direction first. Extend the
    orthonormal list $e_1, e_2, \ldots, e_m$ to an orthonormal basis
    $e_1, \ldots, e_n,$ where $n \geq m.$ Then, we have $v = v_1 e_1
    + \ldots + v_n e_n$ for some $v_k \in \mathbb F.$ Hence,
    $$
    \sum_{i = 1}^{m} |\langle v, e_i\rangle |^2 = \sum_{i = 1}^{m} v_i \bar v_i.
    $$
    By our hypothesis, this is equal to $$
    v_1 \bar v_1 + v_2 \bar v_2 + \ldots + v_n \bar v_n,
    $$
    meaning $v_{m + 1} \bar v_{m + 1} + \ldots + v_{n} \bar v_{n}=
    0.$ Each term is non-negative, meaning we need
    $$
    v_{i} \bar v_i = 0, \ \forall m + 1 \leq i  \leq n,
    $$
    which forces $v_i = 0.$

    Therefore, $v = v_1 e_1 + \ldots + v_m e_m \in \text{span}(e_1,
    \ldots, e_m).$

    Conversely, assume $v \in \text{span}(e_1, \ldots, e_m),$ meaning
    for some $v_j \in \mathbb F,$ $v = v_1 e_1 + \ldots + v_me_m.$
    Then, the inner product $\langle v, e_j \rangle$ resolves to
    $v_j$ since $\langle e_i, e_j \rangle = 0$ whenever $i \neq j.$
    Hence, we have
    $$
    \sum_{i = 1}^{m} | \langle v, e_i \rangle|^2 = \sum_{i = 1}^{m}
    v_i \bar v_i = ||v||^2
    $$
    where $|| \cdot ||$ is the standard norm over $\mathbb F.$

    Therefore,
    \begin{align*}
      \sum_{i = 1}^{m} |\langle v, e_i\rangle |^2 = \sum_{i = 1}^{m}
      v_i \bar v_i.
    \end{align*}
    if and only if $$
    v \in \text{span}(e_1, \ldots, e_m)
    $$
    as required.

  \item By the Riesz Representation Theorem, any linear functional
    $\varphi \in V'$ can be written in the form $\phi(v) = \langle v,
    u \rangle$ for some $u \in V.$

    Consider the dual basis $\{\delta_{i}\}_{i = 1}^{n}$ satisfying
    $\delta_i(v_j) = 1$ if $i =j,$ and $0$ otherwise. By the Riesz
    representation theorem, for every such $\delta_i$ there is some
    $u_i$ such that for all $v \in V,$ $\delta_i(v) = \langle v, u_i
    \rangle,$ meaning $$
    \langle v_j, u_i \rangle = \delta_i(v_j) =
    \begin{cases}
      1 & i = j \\
      0 & \text{otherwise}
    \end{cases}.
    $$

    Thus, it remains to prove that $\{u_j\}_{j = 1}^{n}$ is a basis. Assume $$
    \sum_{i = 1}^{n}\lambda_i u_i = 0
    $$
    for some $\lambda_j \in \mathbb F.$ Then, taking the inner
    product with an arbitrary $v,$ and using conjugate linearity, we have
    $$
    \sum_{i = 1}^{n}\bar \lambda_i\langle v, u_i \rangle = 0
    $$
    for all $v \in V,$ or equivalently,
    $$
    \left(\sum_{i = 1}^{n} \bar \lambda_i \delta_i\right)(v) = 0
    $$
    for all $v,$ implying $$
    \sum_{i = 1}^{n} \bar \lambda_i \delta_i = 0 \in V'.
    $$
    Since $\{\delta_i\}_{i = 1}^{n}$ is a basis, this forces $\bar
    \lambda_i = 0$ for all $i,$ meaning $\lambda_i = 0$ over all $i$,
    and hence $u_j$ is linearly independent. Any linearly independent
    set of size $\dim V$ is a basis over $V,$ meaning $\{u_i\}_{i =
    1}^{n}$ is a basis.

    Hence, there exists a basis $\{u_i\}_{i =1}^{n}$ with $$
    \begin{cases}
      1 & i = j \\
      0 & \text{otherwise}
    \end{cases}$$
    as required. \qed
\end{enumerate}

\subsection{Evaluating Proofs} The structure of the proof is correct,
but you need roots, i.e., $x = (a_1, a_2\sqrt2, a_3 \sqrt 3, \ldots)$
and $b = (a_1, a_2/\sqrt 2, a_2 / \sqrt 3, \ldots).$

\subsection{Using Computers}
\begin{center}
  \includegraphics[width=0.75\linewidth]{math3406/Screenshot
  2026-11-19 at 7.01.58 pm.png}
\end{center}
The claim is correct.
\end{document}
